\chapter{Publier}

Un projet de ce genre produit des articles, des fichiers de données et un site. Tous se
construisent depuis le dépôt, par des commandes qu'on peut relancer.

\section{Des articles reproductibles}

\begin{regle}
Un chiffre cité dans un article se recalcule par une commande, et la commande est dans le
dépôt.
\end{regle}

Les nombres changent pendant que l'article s'écrit : pages lues, théorèmes démontrés,
entrées d'un registre. Si chaque chiffre est recompté par un script, la mise à jour d'une
version révisée prend une minute, et la méthode de comptage est écrite quelque part. Les
figures se dessinent de la même façon, depuis les données, en TikZ, pour que leurs
étiquettes soient composées dans les polices du texte.

\begin{terminal}
# redessiner les figures depuis les données
node figures.mjs
# compiler, sans installation de TeX
tectonic -X compile article.tex
\end{terminal}

L'outil \texttt{tectonic} compile un document \LaTeX{} en téléchargeant les paquets dont il
a besoin ; une seule commande suffit sur n'importe quelle machine, y compris celle de
l'intégration continue.

\section{Des formats d'échange}

Certaines revues demandent un format structuré à côté du PDF : par exemple TEI, un format
XML des humanités numériques, validé contre le schéma de la revue. On écrit un convertisseur
une fois, et la validation entre dans les vérifications automatiques : un fichier qui ne
passe plus le schéma fait échouer la construction.

\section{Un site construit depuis le dépôt}

Le site du projet (une page par document, des registres, une page de relecture) est
construit par un outil de pages statiques et déployé à chaque envoi sur la branche
principale. Les mathématiques s'y affichent avec KaTeX, depuis le même \LaTeX{} que les PDF.
Les documents dont on n'a pas les droits n'y figurent qu'avec l'interrupteur allumé
(chapitre~\ref{ch:droits}).

\section{Soumettre et réviser}

On étiquette la version soumise (chapitre 1). À la révision, \cmd{git diff} depuis
l'étiquette donne la liste des changements, qu'on résume en quelques lignes pour l'éditeur.
On ne soumet jamais deux fois le même article comme s'il était nouveau : une version
révisée passe par la soumission existante.
